\documentclass[11pt]{article}
\usepackage[margin=1in]{geometry}
\usepackage[T1]{fontenc}
\usepackage[utf8]{inputenc}
\usepackage{array}
\usepackage{longtable}
\usepackage{hyperref}
\usepackage{xcolor}
\usepackage{listings}
\lstset{basicstyle=\ttfamily\small,breaklines=true,columns=fullflexible,frame=single}
\setlength{\parindent}{0pt}
\setlength{\parskip}{0.65em}
\title{Lab 2: Pipelining and Superscalar Concepts With SimpleScalar}
\author{Computer Architecture Laboratory}
\date{}
\begin{document}
\maketitle
\section{Objective}
In this lab, students use \texttt{sim-outorder} to explore pipelined and superscalar processor architectures. Students configure and simulate in-order and out-of-order processor models, compare IPC, and use a pipeline trace to reason about data forwarding.
All commands are written for the current repository layout. Run setup commands from the cloned repository root.
\begin{lstlisting}
cd /path/to/SimpleScalar
\end{lstlisting}
\section{Background: The \texttt{sim-outorder} Simulator}
\texttt{sim-outorder} is a detailed microarchitectural timing simulator. It models branch prediction, caches, out-of-order execution, functional units, an instruction window, a load-store queue, and in-order commit.
Get help for the simulator:
\begin{lstlisting}
./sim-outorder -h
\end{lstlisting}
\section{Setup}
Use the PISA build for this lab:
\begin{lstlisting}
make clean
make config-pisa
make
\end{lstlisting}
Create a workspace:
\begin{lstlisting}
mkdir -p labs
\end{lstlisting}
Copy the default configuration and the test program:
\begin{lstlisting}
cp config/default.cfg labs/
cp tests-pisa/bin.little/test-math labs/
\end{lstlisting}
\section{Part 1: In-Order Pipelined Processor}
\subsection{Procedure}
Create \texttt{config\_a.cfg}:
\begin{lstlisting}
cp labs/default.cfg labs/config_a.cfg
\end{lstlisting}
Open \texttt{labs/config\_a.cfg} in a text editor and set:
\begin{lstlisting}
-decode:width 1
-fetch:ifqsize 1
-lsq:size 8
-ruu:size 8
-issue:width 1
-issue:inorder true
-res:ialu 1
-res:imult 1
-res:memport 1
-res:fpalu 1
-res:fpmult 1
\end{lstlisting}
Meaning of the parameters:
\begin{center}
\scriptsize
\begin{longtable}{|p{0.24\linewidth}|p{0.24\linewidth}|}
\hline
Parameter & Meaning \\
\hline
\endfirsthead
\hline
Parameter & Meaning \\
\hline
\endhead
\texttt{-decode:width} & Number of instructions decoded per cycle. \\
\hline
\texttt{-fetch:ifqsize} & Instruction fetch queue size. \\
\hline
\texttt{-lsq:size} & Number of memory operations that can be active at once. \\
\hline
\texttt{-ruu:size} & Register update unit size; the main instruction window. \\
\hline
\texttt{-issue:width} & Number of instructions that can be issued per cycle. \\
\hline
\texttt{-issue:inorder} & Forces program-order issue when set to \texttt{true}. \\
\hline
\texttt{-res:ialu} & Number of integer ALUs. \\
\hline
\texttt{-res:imult} & Number of integer multiply/divide units. \\
\hline
\texttt{-res:memport} & Number of memory ports. \\
\hline
\texttt{-res:fpalu} & Number of floating-point ALUs. \\
\hline
\texttt{-res:fpmult} & Number of floating-point multiply/divide units. \\
\hline
\end{longtable}
\normalsize
\end{center}
Run the simulation:
\begin{lstlisting}
cd labs
../sim-outorder -config config_a.cfg -ptrace config_a.trc 0:1024 -redir:sim sim_configa.out ./test-math
\end{lstlisting}
View the IPC:
\begin{lstlisting}
grep sim_IPC sim_configa.out
\end{lstlisting}
View the pipeline trace:
\begin{lstlisting}
../pipeview.pl config_a.trc | less
\end{lstlisting}
\subsection{Questions for Part 1}
\begin{enumerate}
\item Based on the settings in \texttt{config\_a.cfg}, describe the processor architecture.
\item What is the final value of \texttt{sim\_IPC}?
\item Is data forwarding implemented? To answer, examine the pipeline trace for a RAW hazard. Look for a sequence like:
\end{enumerate}
\begin{lstlisting}
ag: sll  r2,r16,2
ah: addu r3,r3,r2
\end{lstlisting}
Can you find a cycle where \texttt{addu} is in the Execute stage while \texttt{sll} is in the Write-Back stage?
\section{Part 2: Out-of-Order Superscalar Processor}
\subsection{Procedure}
Make sure you are in the \texttt{labs/} directory:
\begin{lstlisting}
cd /path/to/SimpleScalar/labs
\end{lstlisting}
Create \texttt{config\_b.cfg}:
\begin{lstlisting}
cp config_a.cfg config_b.cfg
\end{lstlisting}
Open \texttt{config\_b.cfg} and apply the same modifications as Part 1, with one exception:
\begin{lstlisting}
-issue:inorder false
\end{lstlisting}
Run the simulation:
\begin{lstlisting}
../sim-outorder -config config_b.cfg -redir:sim sim_configb.out ./test-math
\end{lstlisting}
View the IPC:
\begin{lstlisting}
grep sim_IPC sim_configb.out
\end{lstlisting}
\subsection{Questions for Part 2}
\begin{enumerate}
\item Describe the new configuration. What is the key difference compared with Part 1?
\item What is the \texttt{sim\_IPC} for this configuration?
\item Compare the IPC from Part 1 and Part 2. What does this tell you about the benefit of out-of-order execution for this workload?
\item In an out-of-order processor, the Commit stage ensures results are written back in program order. Why is in-order commit necessary?
\end{enumerate}
\end{document}
